\subsection{超参数对照与讨论}
为讨论参数影响，另整理早期探索记录。这部分与后文20局最终榜单分开：表\ref{tab:hyper}各模型均从头训练10万局，训练种子为42，测试种子为20000至20499，共500局；前70\%训练使用$\alpha=0.1$，之后为0.03，前80\%使用$\epsilon=0.02$，之后为0.002。训练局数和参数日程一致，但不同策略的交互步数与耗时仍有差异。

\begin{table}[htbp]
\centering
\caption{历史TD对照：相同训练局数下的参数与结构变化}\label{tab:hyper}
\begin{tabularx}{\linewidth}{@{}Yrrp{0.24\linewidth}@{}}\toprule
设置 & 平均总和 & 达到256 & 相对基准的变化\\\midrule
对称5-tuple，初值0（基准） & 605.76 & 89.0\% & ---\\
5-tuple，不共享对称 & 417.69 & 72.2\% & 对称特征\\
对称4-tuple & 432.88 & 74.0\% & 局部模式结构\\
对称5-tuple，初值300 & 540.78 & 94.4\% & 乐观初始化\\
对称5-tuple，合并分奖励 & 592.64 & 88.4\% & 奖励目标及尺度\\
对称5-tuple，首次256加奖 & 482.92 & 86.8\% & 奖励增加达标项\\
对称6-tuple，首次256加奖 & 401.52 & 66.6\% & 对比上一行的模式\\\bottomrule
\end{tabularx}
\end{table}

\textbf{表示与初始化。}对称共享使平均总和由417.69提高到605.76；它使旋转、镜像局面共享经验，是本轮较明确的改进。4-tuple低于5-tuple，而采用达标奖励的6-tuple也低于同奖励的5-tuple。由于模式数量和表容量同时变化，这说明当前训练量下更大的局部模式未必更好，不能单独归因于特征长度。初值300将达标率从89.0\%提高到94.4\%，平均总和却降低约65，说明更容易达到256与继续累积总和并不总是一致；现有记录没有比较多个非零初值，不能认定300最优。

\textbf{训练量、学习率与探索率。}表\ref{tab:budget}比较10万局与另一次从头训练30万局的结果，测试仍使用同一组500个种子。三组平均总和均上升，但乐观版本的达标率下降，更长训练没有让两项指标同步提高。两轮的退火时间点也随总预算后移，不能视为同一训练轨迹上的两个检查点。

\begin{table}[!htb]
\centering
\caption{训练预算变化：平均总和与达标率}\label{tab:budget}
\begin{tabularx}{\linewidth}{@{}Yrrrr@{}}\toprule
设置 & 10万局总和 & 30万局总和 & 10万局达标 & 30万局达标\\\midrule
普通对称5-tuple & 605.76 & 629.84 & 89.0\% & 93.0\%\\
乐观对称5-tuple & 540.78 & 618.68 & 94.4\% & 90.4\%\\
5-tuple加达标奖励 & 482.92 & 595.00 & 86.8\% & 90.6\%\\\bottomrule
\end{tabularx}
\end{table}

训练日志中降低学习率与探索率后回报上升，但日志包含探索动作，又与训练进度重合，无法分离各因素的贡献。在另一组从同一权重续训20万局的对照中，固定$\alpha=0.01$，$\epsilon=0.0005$与0分别取得708.79和699.23的平均总和；300局配对bootstrap的差值95\%区间为$[-16.4,35.3]$，尚不足以确认保留微量探索的优势。

\textbf{前瞻深度与计算成本。}固定10万局普通5-tuple权重，只改变决策时的前瞻深度，结果见表\ref{tab:depth}。深度3相对深度2提高约6.7\%的总和，耗时约为14.4倍。该结果支持浅层搜索与TD估值结合；计时来自历史编译实现，排除了编译预热，不能代表当前Python入口速度或完整Expectimax的成本。

\begin{table}[!htb]
\centering
\caption{固定TD权重的搜索深度对照（500局）}\label{tab:depth}
\begin{tabularx}{\linewidth}{@{}Yrrr@{}}\toprule
决策方式 & 平均总和 & 达到256 & 每步耗时/ms\\\midrule
直接TD估值（深度1） & 605.76 & 89.0\% & 0.00076\\
TD与深度2搜索 & 670.88 & 94.2\% & 0.00615\\
TD与深度3搜索 & 716.07 & 95.8\% & 0.08855\\\bottomrule
\end{tabularx}
\end{table}

\textbf{结论边界。}现有证据支持对称共享，以及在允许更多决策时间时加深前瞻。乐观初始化和达标奖励需要同时观察总和与达标率；合并分奖励还改变了奖励尺度，不能只解释为目标差异。学习率、MCTS的模拟次数与探索系数、PPO的clip和折扣因子都缺少逐项多取值对照，目前只能说明所用配置及其作用。以上每种配置只有一个训练种子，评测置信区间也不包含重新训练的波动，不能称为全局最优超参数。
